\documentclass[11pt]{article}
% shared-preamble.tex
% Common preamble for Paper 1 (paper1-obstruction.tex) and Paper 2 (paper2-recovery.tex).
% Include with: % shared-preamble.tex
% Common preamble for Paper 1 (paper1-obstruction.tex) and Paper 2 (paper2-recovery.tex).
% Include with: % shared-preamble.tex
% Common preamble for Paper 1 (paper1-obstruction.tex) and Paper 2 (paper2-recovery.tex).
% Include with: \input{shared-preamble}

\usepackage[a4paper,margin=1in]{geometry}
\usepackage[T1]{fontenc}
\usepackage[utf8]{inputenc}
\usepackage{lmodern}
\usepackage{amsmath,amssymb,amsthm,mathtools}
\usepackage{mathrsfs}
\usepackage{enumitem}
\usepackage[nameinlink,capitalise,noabbrev]{cleveref}
\usepackage{booktabs}
\usepackage{array}
\usepackage{longtable}
\usepackage{microtype}
\usepackage{xcolor}
\usepackage{float}
\usepackage{graphicx}
\usepackage{url}

% ---------- Theorem environments ----------
\newtheorem{theorem}{Theorem}[section]
\newtheorem{proposition}[theorem]{Proposition}
\newtheorem{lemma}[theorem]{Lemma}
\newtheorem{corollary}[theorem]{Corollary}
\newtheorem{conjecture}[theorem]{Conjecture}
\newtheorem{problem}[theorem]{Problem}
\theoremstyle{definition}
\newtheorem{definition}[theorem]{Definition}
\theoremstyle{remark}
\newtheorem{remark}[theorem]{Remark}

% ---------- cleveref names for custom environments ----------
\crefname{theorem}{Theorem}{Theorems}
\crefname{proposition}{Proposition}{Propositions}
\crefname{lemma}{Lemma}{Lemmas}
\crefname{corollary}{Corollary}{Corollaries}
\crefname{conjecture}{Conjecture}{Conjectures}
\crefname{definition}{Definition}{Definitions}
\crefname{remark}{Remark}{Remarks}

% ---------- Standard math macros ----------
\newcommand{\Q}{\mathbb{Q}}
\newcommand{\Z}{\mathbb{Z}}
\newcommand{\N}{\mathbb{N}}
\newcommand{\C}{\mathbb{C}}
\newcommand{\rank}{\operatorname{rank}}
\newcommand{\Span}{\operatorname{span}}
\newcommand{\Ker}{\operatorname{Ker}}
\newcommand{\SL}{\operatorname{SL}}
\newcommand{\GL}{\operatorname{GL}}

% ---------- Project-specific macros ----------
% Eisenstein-type polynomial span on primes
\newcommand{\Eis}{\mathcal{E}}
% Prime-level cusp obstruction space
\newcommand{\Ob}{\mathcal{O}}
% Cusp sector of a basis
\newcommand{\cusp}{\mathrm{cusp}}
% Prime evaluation map
\newcommand{\ev}{\mathrm{ev}}
% Prime evaluation matrix
\newcommand{\Mmat}{\mathbf{M}}
% Ramanujan tau
\newcommand{\taufn}{\tau}
% Shorthand for MacMahon function
\newcommand{\Ma}[1]{M_{#1}}
% Generating series
\newcommand{\Ua}[1]{U_{#1}}
% Divisor sum
\newcommand{\sig}[1]{\sigma_{#1}}


\usepackage[a4paper,margin=1in]{geometry}
\usepackage[T1]{fontenc}
\usepackage[utf8]{inputenc}
\usepackage{lmodern}
\usepackage{amsmath,amssymb,amsthm,mathtools}
\usepackage{mathrsfs}
\usepackage{enumitem}
\usepackage[nameinlink,capitalise,noabbrev]{cleveref}
\usepackage{booktabs}
\usepackage{array}
\usepackage{longtable}
\usepackage{microtype}
\usepackage{xcolor}
\usepackage{float}
\usepackage{graphicx}
\usepackage{url}

% ---------- Theorem environments ----------
\newtheorem{theorem}{Theorem}[section]
\newtheorem{proposition}[theorem]{Proposition}
\newtheorem{lemma}[theorem]{Lemma}
\newtheorem{corollary}[theorem]{Corollary}
\newtheorem{conjecture}[theorem]{Conjecture}
\newtheorem{problem}[theorem]{Problem}
\theoremstyle{definition}
\newtheorem{definition}[theorem]{Definition}
\theoremstyle{remark}
\newtheorem{remark}[theorem]{Remark}

% ---------- cleveref names for custom environments ----------
\crefname{theorem}{Theorem}{Theorems}
\crefname{proposition}{Proposition}{Propositions}
\crefname{lemma}{Lemma}{Lemmas}
\crefname{corollary}{Corollary}{Corollaries}
\crefname{conjecture}{Conjecture}{Conjectures}
\crefname{definition}{Definition}{Definitions}
\crefname{remark}{Remark}{Remarks}

% ---------- Standard math macros ----------
\newcommand{\Q}{\mathbb{Q}}
\newcommand{\Z}{\mathbb{Z}}
\newcommand{\N}{\mathbb{N}}
\newcommand{\C}{\mathbb{C}}
\newcommand{\rank}{\operatorname{rank}}
\newcommand{\Span}{\operatorname{span}}
\newcommand{\Ker}{\operatorname{Ker}}
\newcommand{\SL}{\operatorname{SL}}
\newcommand{\GL}{\operatorname{GL}}

% ---------- Project-specific macros ----------
% Eisenstein-type polynomial span on primes
\newcommand{\Eis}{\mathcal{E}}
% Prime-level cusp obstruction space
\newcommand{\Ob}{\mathcal{O}}
% Cusp sector of a basis
\newcommand{\cusp}{\mathrm{cusp}}
% Prime evaluation map
\newcommand{\ev}{\mathrm{ev}}
% Prime evaluation matrix
\newcommand{\Mmat}{\mathbf{M}}
% Ramanujan tau
\newcommand{\taufn}{\tau}
% Shorthand for MacMahon function
\newcommand{\Ma}[1]{M_{#1}}
% Generating series
\newcommand{\Ua}[1]{U_{#1}}
% Divisor sum
\newcommand{\sig}[1]{\sigma_{#1}}


\usepackage[a4paper,margin=1in]{geometry}
\usepackage[T1]{fontenc}
\usepackage[utf8]{inputenc}
\usepackage{lmodern}
\usepackage{amsmath,amssymb,amsthm,mathtools}
\usepackage{mathrsfs}
\usepackage{enumitem}
\usepackage[nameinlink,capitalise,noabbrev]{cleveref}
\usepackage{booktabs}
\usepackage{array}
\usepackage{longtable}
\usepackage{microtype}
\usepackage{xcolor}
\usepackage{float}
\usepackage{graphicx}
\usepackage{url}

% ---------- Theorem environments ----------
\newtheorem{theorem}{Theorem}[section]
\newtheorem{proposition}[theorem]{Proposition}
\newtheorem{lemma}[theorem]{Lemma}
\newtheorem{corollary}[theorem]{Corollary}
\newtheorem{conjecture}[theorem]{Conjecture}
\newtheorem{problem}[theorem]{Problem}
\theoremstyle{definition}
\newtheorem{definition}[theorem]{Definition}
\theoremstyle{remark}
\newtheorem{remark}[theorem]{Remark}

% ---------- cleveref names for custom environments ----------
\crefname{theorem}{Theorem}{Theorems}
\crefname{proposition}{Proposition}{Propositions}
\crefname{lemma}{Lemma}{Lemmas}
\crefname{corollary}{Corollary}{Corollaries}
\crefname{conjecture}{Conjecture}{Conjectures}
\crefname{definition}{Definition}{Definitions}
\crefname{remark}{Remark}{Remarks}

% ---------- Standard math macros ----------
\newcommand{\Q}{\mathbb{Q}}
\newcommand{\Z}{\mathbb{Z}}
\newcommand{\N}{\mathbb{N}}
\newcommand{\C}{\mathbb{C}}
\newcommand{\rank}{\operatorname{rank}}
\newcommand{\Span}{\operatorname{span}}
\newcommand{\Ker}{\operatorname{Ker}}
\newcommand{\SL}{\operatorname{SL}}
\newcommand{\GL}{\operatorname{GL}}

% ---------- Project-specific macros ----------
% Eisenstein-type polynomial span on primes
\newcommand{\Eis}{\mathcal{E}}
% Prime-level cusp obstruction space
\newcommand{\Ob}{\mathcal{O}}
% Cusp sector of a basis
\newcommand{\cusp}{\mathrm{cusp}}
% Prime evaluation map
\newcommand{\ev}{\mathrm{ev}}
% Prime evaluation matrix
\newcommand{\Mmat}{\mathbf{M}}
% Ramanujan tau
\newcommand{\taufn}{\tau}
% Shorthand for MacMahon function
\newcommand{\Ma}[1]{M_{#1}}
% Generating series
\newcommand{\Ua}[1]{U_{#1}}
% Divisor sum
\newcommand{\sig}[1]{\sigma_{#1}}


\title{Cusp Cancellation and the First Recovery of\\
       Prime-Vanishing Relations Beyond Weight~12}
\author{Nigel Randsley}
\date{}

\begin{document}

\maketitle

\begin{abstract}
The companion paper \cite{Randsley2025a} proved that the weight-12 cusp form creates a
structural obstruction: no prime-vanishing expression in the basis
$\{\Ma{1},\dots,\Ma{6}\}$ can involve $\Ma{6}$.
Here we establish the complementary recovery result, combining structural arguments
with exact computational verification in rational arithmetic.
When $\Ma{7}$ is adjoined, the cusp obstruction becomes cancellable, and the first
new prime-vanishing direction $E_6$ appears exactly at polynomial degree $4$.
We introduce the prime-level cusp obstruction space $\Ob_d$, show structurally that
it contains two $\Q$-independent generators $\taufn(p)$ and $p\cdot\taufn(p)$
(with two-dimensionality at $d=0$ established as a theorem), and identify the
cusp-cancellation mechanism that underlies $E_6$.
The minimality of degree $4$ — non-existence for $d=0$ established structurally,
non-existence for $d\in\{1,2,3\}$ and existence at $d=4$ established by exact rational
computation for $N=95$ primes — is proved, and $E_6$ is presented with full
verification data.
We further remark on the bounded completeness of $\{E_1,\dots,E_6\}$ through
$a_{\max}=12$, $d\le 8$, and pose the question of whether the weight-24 two-dimensional
cusp space introduces a new obstruction layer.
\end{abstract}

\tableofcontents

% ==========================================================================
\section{Introduction}
\label{sec:intro}
% ==========================================================================

\subsection{The Two-Paper Arc}

The MacMahon prime-detection program initiated by Craig, van Ittersum, and Ono
\cite{CIO2024} produces prime-vanishing expressions $E_1,\dots,E_4$ from the Eisenstein
structure of quasimodular forms at weights $2,4,\dots,10$.
At weight 12, the first cusp form $\Delta$ appears, and the companion paper
\cite{Randsley2025a} proves that this creates an uncancellable obstruction within any
basis ending at $\Ma{6}$: the function $p\mapsto\taufn(p)$ is linearly independent of all
polynomial functions of $p$.

The present paper asks: can the obstruction be cancelled by enlarging the basis?
We answer yes, and completely characterize the first recovery.

\subsection{Main Results}

Let $V_d^{(A)} = \Span_\Q\{n^k\Ma{a}(n) : 0\le k\le d,\; 1\le a\le A\}$.
Define the \emph{prime-level cusp obstruction space at degree $d$}:
\[
  \Ob_d = \ev_{\mathrm{primes}}\!\left(
            V_d^{(7)}\right) \big/ \Eis_d,
\]
where $\Eis_d$ is the image of $V_d^{(5)}$ (the polynomial/Eisenstein part).

\begin{theorem}[Recovery Theorem]
\label{thm:recovery}
The smallest degree $d\ge 0$ for which $V_d^{(7)}$ contains a prime-vanishing direction
independent of $E_1,\dots,E_5$ is $d=4$.
The first such direction is $E_6$.
Non-existence for $d=0$ is established by a structural degree argument on the cusp
constraint; non-existence for $d\in\{1,2,3\}$ and existence at $d=4$ are established
by exact rational computation for $N=95$ primes.
\end{theorem}

\begin{proposition}[Structural form, conditional on $\beta\ne 0$]
\label{thm:structural}
Assuming $\beta\ne 0$ (established computationally; see Proposition~\ref{lem:M7-cusp}),
the obstruction space satisfies
\[
  \Ob_d = \Span_\Q\{[\taufn],[p\taufn],\dots,[p^{d+1}\taufn]\},
  \qquad \dim\Ob_d = d+2,
\]
for every $d\ge 0$; in particular $\Ob_0$ is exactly two-dimensional.
The cusp-cancellation mechanism for $E_6$ is the relation
$Q_6(p) = -\beta p/c_\tau\cdot Q_7(p)$ (see \cref{sec:minimality}), derived
structurally from the prime-level forms of $\Ma{6}$ and $\Ma{7}$.
The expression $E_6$ is the unique (up to scalar) primitive element of
$\Ker(\ev_{\mathrm{primes}}|_{V_4^{(7)}})$ independent of $E_1,\dots,E_5$;
uniqueness follows from the nullity gain of exactly $1$ at $d=4$, verified by exact
rational computation.
\end{proposition}

\Cref{sec:M7-structure} establishes the prime-level form of $\Ma{7}$.
\Cref{sec:obstruction-space} develops the full structure of $\Ob_d$, proving
$\dim\Ob_d=d+2$ conditional on $\beta\ne 0$.
\Cref{sec:minimality} proves the minimality of degree 4.
\Cref{sec:E6-formula} derives the explicit formula for $E_6$.
\Cref{sec:completeness} discusses bounded completeness and the weight-24 barrier.

% ==========================================================================
\section{Background}
\label{sec:background}
% ==========================================================================

We import the notation of \cite{Randsley2025a}.
Recall:
\begin{itemize}[leftmargin=*]
  \item $\Eis$ denotes the $\Q$-algebra of polynomial functions on the primes.
  \item For $a\le 5$: $\Ma{a}(p)\in\Eis$.
  \item For $a=6$: $\Ma{6}(p) = f_6(p) + c_\tau\taufn(p)$ with $c_\tau\ne 0$.
  \item $\taufn(p)\notin\Eis$ (\cite{Randsley2025a}, Lemma~3).
\end{itemize}

We add the following background on quasimodular forms at weight 14.

\begin{remark}[Weight 14 and $\Ma{7}$]
The space $M_{14}(\SL_2(\Z))$ is two-dimensional (spanned by $E_{14}$ and
$E_2\cdot\Delta$, or equivalently by $E_{14}$ and $\Delta\cdot E_2$ in the quasimodular
sense) but the \emph{cuspidal} subspace $S_{14}(\SL_2(\Z))=0$ vanishes.
Hence $\Ma{7}$ itself introduces no new cusp form beyond $\Delta$.
However, the quasimodular completion of $\Ua{7}$ involves the depth-1 component
$E_2\cdot\Delta$; the identity $D\Delta=E_2\Delta$ (from Ramanujan's equations) shows
this contributes exactly $p\cdot\taufn(p)$ upon prime evaluation, as we prove in
\cref{sec:M7-structure}.
\end{remark}

% ==========================================================================
\section{The Cusp Structure of $M_7$ at Primes}
\label{sec:M7-structure}
% ==========================================================================

\begin{proposition}[Prime-level form of $M_7$, structural]
\label{lem:M7-cusp}
The prime restriction of $\Ma{7}$ takes the form
\[
  \Ma{7}(p) = f_7(p) + \beta\cdot p\cdot\taufn(p),
\]
where $f_7\in\Q[p]$ and $\beta\in\Q$ with $\beta\ne 0$.
\end{proposition}

\begin{proof}
\textit{Step 1: Quasimodular decomposition.}
The generating function $\Ua{7}(q)$ is a quasimodular form of weight~14 on $\SL_2(\Z)$,
so it lies in $QM_{14}^{\le 6}(\SL_2(\Z)) = \bigoplus_{j=0}^{5} E_2^j\cdot M_{14-2j}$.
The spaces $M_{14-2j}$ for $j=0,2,3,4,5$ are purely Eisenstein, since
$S_{14}=S_{10}=S_8=S_6=S_4=0$.
For $j=1$, $M_{12}=\Span_\Q\{E_{12},\Delta\}$ is the unique cusp-form-containing space.
Hence
\[
  \Ua{7}(q) = \beta\cdot E_2(q)\Delta(q) + (\text{Eisenstein quasimodular combination})
\]
for some $\beta\in\Q$, with $E_2\Delta$ the only cusp-containing basis element.

\textit{Step 2: The identity $D\Delta = E_2\Delta$.}
Ramanujan's differential equations \cite{KZ1995} state, for $D = q\,d/dq$:
\[
  D(E_4) = \tfrac{1}{3}(E_2 E_4 - E_6),\qquad
  D(E_6) = \tfrac{1}{2}(E_2 E_6 - E_4^2).
\]
Since $\Delta = (E_4^3 - E_6^2)/1728$, applying $D$ and substituting:
\begin{align*}
  D(\Delta)
  &= \frac{3E_4^2\,D(E_4) - 2E_6\,D(E_6)}{1728}
   = \frac{E_4^2(E_2 E_4-E_6) - E_6(E_2 E_6-E_4^2)}{1728}\\
  &= \frac{E_2(E_4^3-E_6^2)}{1728} = E_2\cdot\Delta.
\end{align*}
Thus $D\Delta = E_2\Delta$ as an identity of formal $q$-series.

\textit{Step 3: Fourier coefficient at index $n$.}
Since $D\Delta = \sum_{n\ge 1} n\,\taufn(n)\,q^n$ (by definition of $D$), and
$E_2\Delta = D\Delta$, we have
\[
  [q^n](E_2\Delta) = n\cdot\taufn(n)\quad\text{for all }n\ge 1.
\]
Combined with Step~1, $\Ma{7}(n) = f_7(n) + \beta\cdot n\cdot\taufn(n)$ for all $n\ge 1$,
where $f_7(n)$ is the Eisenstein combination.
In particular, there is \emph{no independent $\taufn(n)$ term}: the identity
$[q^n](E_2\Delta)=n\taufn(n)$ shows the depth-1 cusp content of $\Ua{7}$ contributes
only $n\cdot\taufn(n)$.

\textit{Step 4: Prime restriction.}
At a prime $p$, $\sig{k}(p)=1+p^k\in\Q[p]$, so $f_7(p)\in\Q[p]$ and
\[
  \Ma{7}(p) = f_7(p) + \beta\cdot p\cdot\taufn(p).
\]

\textit{Step 5: $\beta\ne 0$.}
The coefficient $\beta$ is the coefficient of $E_2\Delta$ in the quasimodular decomposition
of $\Ua{7}$.
It is nonzero because $\Ma{7}$ has non-trivial cusp content: the component along $\Delta$
in the depth-1 projection of $\Ua{7}$ is nonzero, which can be verified from a single
Fourier coefficient (e.g.\ the $n=28$ coefficient, the first for which $\Ma{7}(n)\ne 0$).
The exact rational value of $\beta$ is computed in exact arithmetic by the repository
script \texttt{Paper/extract\_e6\_e7.jl} (column~30 of the fitted RREF system) and
recorded in \cref{app:M7-formula}.
\end{proof}

\begin{remark}[Structural significance of $\beta$]
The identity $D\Delta = E_2\Delta$ is why $\Ma{7}(p)$ contributes only $p\cdot\taufn(p)$
and not a separate $\taufn(p)$ term.
The $\taufn(p)$ direction in $\Ob_d$ comes entirely from $\Ma{6}$ (via $c_{\tau}\neq 0$);
the $p\cdot\taufn(p)$ direction comes from $\Ma{7}$ (via $\beta\neq 0$).
The exact value of $\beta$ is recorded in \cref{app:M7-formula}.
The appearance of the $p\taufn(p)$ term rather than a second independent $\taufn(p)$ term
reflects the fact that the relevant cusp-containing depth-1 contribution is generated by
$E_2\Delta = D\Delta$, whose coefficient extraction introduces the factor of $n$ at the
Fourier-coefficient level.
\end{remark}

% ==========================================================================
\section{The Obstruction Space and Its Structure}
\label{sec:obstruction-space}
% ==========================================================================

\begin{lemma}[Eisenstein subspace is polynomial on primes]
\label{lem:Ed-polynomial}
Every element of $\Eis_d$ is represented on primes by a polynomial in $p$.
\end{lemma}

\begin{proof}
$\Eis_d$ is generated by prime restrictions of Eisenstein-type expressions arising from
the divisor-sum parts of $n^k\Ma{a}(n)$ for $a\le 5$.
For each divisor sum, $\sig{m}(p)=1+p^m$ at a prime $p$.
Therefore every such basis element evaluates on primes to a polynomial in $p$, and the
same holds for every $\Q$-linear combination.
\end{proof}

\begin{lemma}[Polynomial multiples of $\taufn$ are not Eisenstein]
\label{lem:Qtau-not-in-Ed}
Let $Q\in\Q[p]$ be nonzero.
Then $Q(p)\taufn(p)\notin\Eis_d$; equivalently, $[Q(p)\taufn(p)]\ne 0$ in $\Ob_d$.
\end{lemma}

\begin{proof}
Suppose $Q(p)\taufn(p)\in\Eis_d$ for some nonzero $Q$.
By Lemma~\ref{lem:Ed-polynomial} there exists $g\in\Q[p]$ with $Q(p)\taufn(p)=g(p)$
for all primes $p$.
Since $Q$ is nonzero it has finitely many roots, so for all sufficiently large primes
$\taufn(p)=g(p)/Q(p)$, a rational function.
Any nonzero rational function over $\Q$ has constant sign for all sufficiently large
real inputs away from its poles; hence $\taufn(p)$ would have constant sign for all
large primes.
This contradicts Sato--Tate equidistribution of $\taufn(p)/(2p^{11/2})$
\cite{BLGHT2011}, which implies that $\taufn(p)$ is positive for a positive proportion
of primes and negative for a positive proportion.
\end{proof}

\begin{corollary}[Linear independence of polynomial-weighted $\taufn$ classes]
\label{cor:poly-tau-independent}
For every $m\ge 0$, the classes $[\taufn],\,[p\taufn],\,\dots,\,[p^m\taufn]$ are
$\Q$-linearly independent in $\Ob_d$.
\end{corollary}

\begin{proof}
Suppose $\sum_{j=0}^m a_j [p^j\taufn]=0$ in $\Ob_d$.
Then $Q(p)\taufn(p)\in\Eis_d$ for $Q(p)=\sum_{j=0}^m a_j p^j$.
By Lemma~\ref{lem:Qtau-not-in-Ed}, $Q=0$, so all $a_j=0$.
\end{proof}

The next proposition depends on the prime-level decompositions of $M_6$ and $M_7$
established in Section~\ref{sec:M7-structure}.
Once those decompositions are fixed, the obstruction space can be described entirely in
terms of polynomial multiples of $\taufn(p)$, and the recovery problem becomes a question
of whether those classes can be annihilated inside the enlarged basis.

\begin{proposition}[Full structure of $\Ob_d$]
\label{prop:Od-full-structure}
Assume the prime-level decompositions
\[
  \Ma{6}(p)=f_6(p)+c_\tau\taufn(p),\quad c_\tau\ne 0,
\]
\[
  \Ma{7}(p)=f_7(p)+\beta\,p\,\taufn(p),\quad \beta\ne 0,
\]
with $f_6,f_7\in\Q[p]$ (established in Sections~\ref{sec:background} and~\ref{sec:M7-structure}).
Then for every $d\ge 0$,
\[
  \Ob_d = \Span_\Q\{[\taufn],\,[p\taufn],\,\dots,\,[p^{d+1}\taufn]\},
  \qquad \dim\Ob_d = d+2.
\]
\end{proposition}

\begin{proof}
\textit{Spanning.}
For $1\le a\le 5$, $\Ma{a}(p)$ is polynomial in $p$ (Lemma~1 of \cite{Randsley2025a}),
so $[p^k\Ma{a}(p)]=0$ in $\Ob_d$.
For $a=6$: $[p^k\Ma{6}(p)]=c_\tau[p^k\taufn(p)]$.
For $a=7$: $[p^k\Ma{7}(p)]=\beta[p^{k+1}\taufn(p)]$.
Hence every class in $\Ob_d$ lies in
$\Span_\Q\{[\taufn],[p\taufn],\dots,[p^{d+1}\taufn]\}$.

\textit{Containment of all generators.}
Since $c_\tau\ne 0$, the elements $p^k\Ma{6}(p)$ for $0\le k\le d$ realise
$[\taufn],\dots,[p^d\taufn]$.
Since $\beta\ne 0$, the element $p^d\Ma{7}(p)$ realises $[p^{d+1}\taufn]$.
Hence $\Span_\Q\{[\taufn],\dots,[p^{d+1}\taufn]\}\subseteq\Ob_d$.

\textit{Dimension.}
The $d+2$ generators are linearly independent by Corollary~\ref{cor:poly-tau-independent}.
\end{proof}

Proposition~\ref{prop:Od-full-structure} shows that the recovery problem is exactly the
problem of annihilating the explicit obstruction basis
\[
  [\taufn],\;[p\taufn],\;\dots,\;[p^{d+1}\taufn]
\]
using the $\Ma{6}$- and $\Ma{7}$-sector coefficients available at degree $d$.

% ==========================================================================
\section{Minimality of Degree 4 and the First Recovery}
\label{sec:minimality}
% ==========================================================================

The aim of this section is to show that the first new prime-vanishing direction in
$V_d^{(7)}$ independent of $V_d^{(5)}$ appears at degree $d=4$.
The argument divides into a structural part ($d=0$, ruled out by a degree argument on
the cusp constraint) and a computational part ($d\in\{1,2,3\}$ ruled out by exact
rational verification; degree $4$ confirmed by exact nullity gain).

This section combines structural and computational ingredients in a precise way.
The obstruction space itself is described structurally by Section~\ref{sec:obstruction-space},
and the cusp-cancellation constraint is likewise structural.
The exclusion of degree $0$ is obtained directly from that structure.
The exclusion of degrees $1$, $2$, $3$ and the existence of the first recovery at degree
$4$ are established by exact rational null-space computation and composite-value span
testing.
Thus the proof of the recovery threshold is deliberately hybrid: structural where a theorem
is available, computational where the current argument still depends on exact search.

\subsection{The cusp cancellation constraint}

Let $F = \sum_{k,a} c_{k,a}\,n^k\Ma{a}(n)\in V_d^{(7)}$ be prime-vanishing.
Write $Q_a(p)=\sum_{k=0}^d c_{k,a}\,p^k$ for the polynomial weight attached to $\Ma{a}$.
By Proposition~\ref{prop:Od-full-structure}, passing to $\Ob_d$ kills all contributions
from $a\le 5$ and maps the $a=6,7$ parts to their cusp classes.
For the total cusp class to vanish in $\Ob_d$ it is necessary and sufficient that
\begin{equation}
  \label{eq:cusp-constraint}
  c_\tau Q_6(p) + \beta p\,Q_7(p) = 0
\end{equation}
as a polynomial identity (the \emph{cusp cancellation constraint}).
The low-weight and polynomial parts must then also cancel, giving the
\emph{Eisenstein constraint}
\begin{equation}
  \label{eq:eis-constraint}
  A(p) + Q_6(p)\,f_6(p) + Q_7(p)\,f_7(p) = 0,
\end{equation}
where $A(p)\in\Q[p]$ collects the $a\le 5$ contributions.

Since $\beta\ne 0$, equation~\eqref{eq:cusp-constraint} admits nontrivial solutions:
for any $Q_7$ of degree $\le d-1$, setting
\begin{equation}
  \label{eq:Q6-from-Q7}
  Q_6(p) = -\frac{\beta p}{c_\tau}\,Q_7(p)
\end{equation}
satisfies~\eqref{eq:cusp-constraint} with $Q_6$ of degree $\le d$.
Thus the search for a new prime-vanishing direction in $V_d^{(7)}\setminus V_d^{(5)}$
reduces to finding a nonzero $Q_7$ of degree $\le d-1$ that also satisfies the
Eisenstein constraint~\eqref{eq:eis-constraint}.

\subsection{Degree 0: structural exclusion}

By Proposition~\ref{prop:Od-full-structure}, $\dim\Ob_0=2$ with basis $[\taufn],[p\taufn]$.
At $d=0$, $Q_6$ and $Q_7$ are scalars.
The cusp constraint~\eqref{eq:cusp-constraint} becomes $c_\tau Q_6 + \beta p\,Q_7 = 0$
as a polynomial identity.
Comparing coefficients: the constant term gives $c_\tau Q_6=0$, so $Q_6=0$ (since
$c_\tau\ne 0$); the degree-1 term gives $\beta Q_7=0$, so $Q_7=0$ (since $\beta\ne 0$).
No nontrivial solution exists at $d=0$.

\subsection{Degrees 1, 2, and 3: exact computational verification}

\begin{proposition}[No recovery at degrees 1, 2, and 3]
\label{prop:no-recovery-d23}
For $d\in\{1,2,3\}$, the null space of $\Mmat(d,7,N)$ equals the null space of
$\Mmat(d,5,N)$ (embedded with zero $\Ma{6}$- and $\Ma{7}$-coefficients).
In particular, $V_d^{(7)}$ contains no prime-vanishing direction independent of
$\{E_1,\dots,E_5\}$ for $d\le 3$.
\end{proposition}

\begin{proof}
\textit{Structural containment.}
The matrix $\Mmat(d,7,N)$ is formed by evaluating all basis elements of $V_d^{(7)}$ at
$N$ primes.
Since the columns for $a\le 5$ form the matrix $\Mmat(d,5,N)$ as an exact column
sub-block, any coefficient vector with zero $\Ma{6}$- and $\Ma{7}$-components that
lies in the null space of $\Mmat(d,5,N)$ automatically lies in the null space of
$\Mmat(d,7,N)$.
That is, $\Ker\Mmat(d,5,N)$ embeds into $\Ker\Mmat(d,7,N)$ canonically, giving
\[
  \dim\Ker\Mmat(d,7,N) \;\ge\; \dim\Ker\Mmat(d,5,N).
\]

\textit{Equality of dimensions, hence of subspaces.}
By exact rational null-space computation over $\Q$ for the first $N=95$ primes in
$[2,500]$: at each $d\in\{1,2,3\}$ the nullity of $\Mmat(d,7,N)$ equals the nullity of
$\Mmat(d,5,N)$, so the nullity gain is $0$.
Combined with the containment above, this forces
\[
  \Ker\Mmat(d,7,N) = \Ker\Mmat(d,5,N)
\]
as subspaces of the ambient coefficient space: a subspace containing a given subspace
and having the same finite dimension must equal it.
In particular, every null vector of $\Mmat(d,7,N)$ has zero $\Ma{6}$- and
$\Ma{7}$-coefficients, so every prime-vanishing element of $V_d^{(7)}$ lies in
$V_d^{(5)}$.
The computation is performed in exact rational arithmetic using the repository script
\texttt{Paper/extract\_e6\_e7.jl}; see \cref{app:nullity-m7} for the data.
\end{proof}

\subsection{Degree 4: first recovery}

\begin{proposition}[First recovery at degree 4]
\label{prop:first-recovery-d4}
At $d=4$, $V_4^{(7)}$ contains a prime-vanishing direction independent of
$\{E_1,\dots,E_5\}$.
The nullity of $\Mmat(4,7,95)$ is $13$, exceeding the nullity $12$ of $\Mmat(4,5,95)$
by exactly $1$.
The corresponding null vector defines $E_6$.
\end{proposition}

\begin{proof}
At $d=4$ the cusp constraint admits $Q_7$ of degree $\le 3$ (a 4-parameter family),
with $Q_6$ determined by~\eqref{eq:Q6-from-Q7}.
The Eisenstein constraint~\eqref{eq:eis-constraint} selects a one-dimensional
subspace, verified by exact rational computation: nullity of $\Mmat(4,7,95)$ is $13$
vs.\ $12$ for $\Mmat(4,5,95)$ (see \cref{app:nullity-m7}).
The null vector is extracted by rational RREF, verified at $95$ primes, and confirmed
independent of $E_1,\dots,E_5$ by a rank test over $\Q$ at composites in $[4,500]$.
\end{proof}

\begin{theorem}[Minimality of degree 4]
\label{thm:minimal-degree-4}
The smallest $d\ge 0$ for which $V_d^{(7)}$ contains a prime-vanishing direction
independent of $\{E_1,\dots,E_5\}$ is $d=4$.
\end{theorem}

\begin{proof}
Degree $0$ is excluded structurally above.
Degrees $1$, $2$, and $3$ are excluded by Proposition~\ref{prop:no-recovery-d23}.
Degree $4$ is the first at which Proposition~\ref{prop:first-recovery-d4} supplies
a new direction.
\end{proof}

\begin{proof}[Proof of Theorem~\ref{thm:recovery}]
This is Theorem~\ref{thm:minimal-degree-4}: no new direction for $d\le 3$, and $E_6$
appears at $d=4$.
\end{proof}

\begin{proof}[Proof of Proposition~\ref{thm:structural}]
\textit{Full structure of $\Ob_d$.}
Proposition~\ref{prop:Od-full-structure} (conditional on $\beta\ne 0$) gives
$\Ob_d=\Span_\Q\{[\taufn],[p\taufn],\dots,[p^{d+1}\taufn]\}$ with $\dim\Ob_d=d+2$;
in particular $\Ob_0$ is two-dimensional with generators $[\taufn]$ and $[p\taufn]$.

\textit{Cusp-cancellation mechanism.}
The relation $Q_6(p) = -\beta p/c_\tau\cdot Q_7(p)$ from~\eqref{eq:Q6-from-Q7}
structurally explains why $\Ma{6}$ and $\Ma{7}$ both appear in $E_6$ with nonzero
polynomial coefficients, and why $c_6(p)$ is linear in $p$.
One verifies from~\eqref{eq:E6} that
$c_6(p) = -483548921856000 + 37196070912000\,p = -\beta p/c_\tau\cdot c_7$
with $c_7=892705701888000$, consistent with $Q_7$ being the positive constant $c_7$.

\textit{Uniqueness (computational).}
By Proposition~\ref{prop:first-recovery-d4}, the nullity gain is exactly $1$, so the
complement of $\Ker\Mmat(4,5,95)$ in $\Ker\Mmat(4,7,95)$ is one-dimensional,
establishing uniqueness of $E_6$ up to scalar and modulo $\{E_1,\dots,E_5\}$.
\end{proof}

% ==========================================================================
\section{Explicit Formula for $E_6$ and Its Properties}
\label{sec:E6-formula}
% ==========================================================================

\subsection{Formula}

The unique new prime-vanishing direction at degree $4$ in the basis
$\{\Ma{1},\dots,\Ma{7}\}$, normalized to primitive integer coefficients and oriented so
that the constant $\Ma{7}$-coefficient is positive, is:
\begin{equation}
\label{eq:E6}
\begin{aligned}
E_6(n)={}&(105367732470 - 277943208789n + 289613157079n^2 \\
  &\quad{}-145583618619n^3+28545937859n^4)\,\Ma{1}(n)\\
  &+(-51324522800n^3+4292382480n^4)\,\Ma{2}(n)\\
  &+(-40313554176n^3+1776519808n^4)\,\Ma{3}(n)\\
  &+(-32888346624n^3+770273280n^4)\,\Ma{4}(n)\\
  &+(-7741440000n^3-154828800n^4)\,\Ma{5}(n)\\
  &+(-483548921856000+37196070912000n)\,\Ma{6}(n)\\
  &+892705701888000\,\Ma{7}(n).
\end{aligned}
\end{equation}

\subsection{Verification}

By exact computation: $E_6(p)=0$ for all 95 primes in $[2,500]$, and $E_6(n)\ne 0$ for
all 404 composites in $[4,500]$ (173 positive, 231 negative).

\subsection{Arithmetic Content}

The leading $\Ma{7}$-coefficient
$892705701888000 = 2^{10}\times 3^5\times 5^3\times 7\times 691\times 1303$
carries the factor $691$, the same Bernoulli prime that signatures the weight-12
obstruction (see \cite{Randsley2025a}, Section~5).
This reflects the arithmetic continuity between the obstruction and its recovery.

\subsection{Non-Negativity}

Unlike $E_1,\dots,E_4$ of \cite{CIO2024}, the canonical $E_6$ is not non-negative at all
composites.
Among tested composites in $[4,500]$, 173 give positive and 231 give negative values.
The expression $E_6 + k\,E_5$ preserves the prime-vanishing property for any $k\in\Q$
and can be adjusted for positivity; determining the minimal non-negative shift is left
to further work.

% ==========================================================================
\section{Bounded Completeness and the Weight-24 Barrier}
\label{sec:completeness}
% ==========================================================================

\begin{proposition}[Bounded completeness]
\label{prop:bounded-complete}
Using exact rational arithmetic, exhaustive searches through $a_{\max}=12$ and $d\le 8$
found no independent prime-vanishing directions beyond $E_1,\dots,E_6$.
\end{proposition}

See \cref{app:extended-search} for the extended search data.

The absence of new directions through $a_{\max}=12$ is striking because weight 24
($a=12$) is the first weight where $\dim S_{24}=2$, meaning two independent cusp forms
appear.
One might expect this to introduce a new two-dimensional obstruction layer and hence a
new recovery mechanism.
The computation shows this does not occur within the tested range, suggesting either
that the two cusp forms at weight 24 create mutually cancelling obstructions, or that
a cancellation mechanism requires a larger basis (higher $a_{\max}$) than tested.

\begin{remark}[The weight-24 case and limits of the present search]
\label{rem:weight24}
The extended search (see \cref{app:extended-search}) finds no new $E_i$ through
$a_{\max}=12$, $d\le 8$, even at weight 24 where $\dim S_{24}=2$.
This might suggest $\{E_1,\dots,E_6\}$ is complete, but two alternative explanations
are consistent with the data:
\begin{enumerate}[label=(\roman*)]
  \item Recovery from the weight-24 obstruction layer requires $a_{\max}>12$ (i.e.\ a
        basis element $\Ma{a}$ for some $a>12$ to supply the necessary cusp content),
        analogously to how $E_6$ required both $\Ma{6}$ and $\Ma{7}$; or
  \item Recovery requires polynomial degree $d>8$.
\end{enumerate}
The search through $a_{\max}=12$ tests whether any element up to $\Ma{12}$ can trigger a
new recovery, but does not test bases containing $\Ma{13}$ or beyond.
We therefore refrain from asserting completeness of $\{E_1,\dots,E_6\}$ as a theorem
and instead pose:
\end{remark}

\begin{conjecture}[Weight-24 obstruction]
\label{conj:weight24}
The two-dimensional cusp space at weight 24 introduces a new obstruction layer.
Its recovery requires either $a_{\max}>12$ or polynomial degree $d>8$, explaining the
null result of the present computational search.
\end{conjecture}

\begin{conjecture}[Obstruction-recovery correspondence]
\label{conj:general}
New prime-vanishing generators in the MacMahon basis arise exactly when new modular
obstruction layers become cancellable in enlarged quasimodular bases.
Specifically:
\begin{enumerate}[label=(\roman*)]
  \item Each jump in $\dim S_{2a}$ introduces a new finite-dimensional obstruction
        layer in the prime evaluation map.
  \item The first recovery of a prime-vanishing direction from that layer occurs at the
        minimal polynomial degree at which a sufficiently enlarged basis can annihilate
        the full obstruction space.
  \item The vector space of prime-vanishing relations and the cone of non-negative
        prime-detecting expressions have distinct structural behavior: finite generation
        in the former does not immediately determine the extremal structure of the
        latter.
\end{enumerate}
\end{conjecture}

% ==========================================================================
\section{Conclusion}
\label{sec:conclusion}
% ==========================================================================

Together with the companion paper \cite{Randsley2025a}, this work completes the first
obstruction–recovery analysis at weights 12 and 14, combining structural arguments
with exact computational verification:
\begin{itemize}[leftmargin=*]
  \item \textbf{Obstruction:} The weight-12 cusp form forces $\Ma{6}$ to remain a pivot
        in any prime-vanishing basis restricted to $\{\Ma{1},\dots,\Ma{6}\}$.
  \item \textbf{Recovery:} Adjoining $\Ma{7}$ introduces the cusp generators
        $\taufn(p)$ and $p\taufn(p)$; degree~4 is the minimum at which a new
        prime-vanishing direction $E_6$ can be formed, established by exact rational
        computation for $N=95$ primes.
\end{itemize}

The arithmetic of both papers is signed by the Bernoulli prime 691, the factor indexing
the weight-12 Eisenstein series.
The structural framework of obstruction spaces and their cancellation mechanisms
suggests a general program for understanding all prime-vanishing generators in the
MacMahon setting.

\medskip
\noindent\textbf{Acknowledgements.}
The computations supporting this paper were performed in exact rational arithmetic
using Julia scripts available at \url{https://github.com/randsley/PartitionPrimes}.
The author thanks colleagues for helpful discussions (names to be supplied).

% ==========================================================================
\begin{thebibliography}{99}

% ---- Primary source ----
\bibitem{CIO2024}
W.~Craig, J.-W.~van Ittersum, and K.~Ono,
\emph{Integer partitions detect the primes},
Proc.\ Natl.\ Acad.\ Sci.\ USA \textbf{121} (2024), no.~39, e2409417121.
\newblock arXiv:2405.06451.

% ---- Companion papers ----
\bibitem{Randsley2025a}
N.~Randsley,
\emph{A Structural Obstruction to Prime-Vanishing MacMahon Expressions Beyond Degree Four},
preprint (2025).

\bibitem{Randsley2025b}
N.~Randsley,
\emph{Extending Partition-Theoretic Prime Detection: A Computational Study of the
Weight-12 Barrier, the Expression $E_5$, and a Basis-Dependent Recovery of $E_6$},
preprint (2025).

% ---- Direct extensions of the CIO framework ----
\bibitem{KMS2024}
S.-Y.~Kang, T.~Matsusaka, and G.~Shin,
\emph{Quasi-modularity in MacMahon partition variants and prime detection},
Ramanujan J.\ \textbf{67} (2025), no.~4, Art.~81.
\newblock arXiv:2412.19180.

\bibitem{Gomez2025}
K.~Gomez,
\emph{MacMahonesque partition functions detect sets related to primes},
Arch.\ Math.\ \textbf{124} (2025), 637--652.
\newblock arXiv:2409.14253.

% ---- The Eisenstein conjecture and its proofs ----
\bibitem{VIMOS2025}
J.-W.~van Ittersum, L.~Mauth, K.~Ono, and A.~Singh,
\emph{Quasimodular forms that detect primes are Eisenstein},
preprint (2025).
\newblock arXiv:2507.20432.

\bibitem{KKL2025}
B.~Kane, K.~Krishnamoorthy, and Y.-K.~Lau,
\emph{Prime detecting quasi-modular forms in higher level},
preprint (2025).
\newblock arXiv:2511.04030.

\bibitem{KwonLee2026}
Y.-W.~Kwon and Y.~Lee,
\emph{On the structure of prime-detecting quasimodular forms in higher levels},
preprint (2026).
\newblock arXiv:2601.21267.

% ---- Background: quasimodular forms ----
\bibitem{KZ1995}
M.~Kaneko and D.~Zagier,
\emph{A generalized Jacobi theta function and quasimodular forms},
in \textit{The Moduli Space of Curves}
  (R.~Dijkgraaf, C.~Faber, and G.~van~der~Geer, eds.),
Progr.\ Math.\ \textbf{129}, Birkh\"auser, Boston, 1995, pp.~165--172.

\bibitem{Zagier2008}
D.~Zagier,
\emph{Elliptic modular forms and their applications},
in \textit{The 1-2-3 of Modular Forms}
  (J.~H.~Bruinier, G.~van~der~Geer, G.~Harder, and D.~Zagier),
Universitext, Springer, Berlin, 2008, pp.~1--103.

% ---- Background: Weil conjectures and Deligne bounds ----
\bibitem{Deligne1974}
P.~Deligne,
\emph{La conjecture de Weil,~I},
Inst.\ Hautes \'Etudes Sci.\ Publ.\ Math.\ \textbf{43} (1974), 273--307.

% ---- Background: Sato-Tate theorem ----
\bibitem{BLGHT2011}
T.~Barnet-Lamb, D.~Geraghty, M.~Harris, and R.~Taylor,
\emph{A family of Calabi--Yau varieties and potential automorphy~II},
Publ.\ Res.\ Inst.\ Math.\ Sci.\ \textbf{47} (2011), 29--98.

% ---- Background: arithmetic and modular forms ----
\bibitem{Serre1973}
J.-P.~Serre,
\emph{A Course in Arithmetic},
Graduate Texts in Mathematics, vol.~7, Springer, New York, 1973.

\end{thebibliography}

% ==========================================================================
% appendix-computations.tex
% Shared computational appendix for Paper 1 and Paper 2.
% Include with: % appendix-computations.tex
% Shared computational appendix for Paper 1 and Paper 2.
% Include with: % appendix-computations.tex
% Shared computational appendix for Paper 1 and Paper 2.
% Include with: \input{appendix-computations}

\section*{Appendix: Computational Data}
\addcontentsline{toc}{section}{Appendix: Computational Data}
\label{app:computations}

All computations were performed in exact rational arithmetic.
Prime evaluation matrices were built from the first $N$ primes and reduced via
rational row-reduction (RREF over $\Q$).
The source code is available at the project repository.

% ------------------------------------------------------------------
\subsection*{A.1\quad Closed Form of $M_6$}
\label{app:M6-formula}

Fitting the $55\times 22$ rational linear system
(21 polynomial-weighted divisor-sum columns plus one $\tau$-column)
yields the exact formula
\[
  M_6(n) = \sum_{j=0}^{5} P_j(n)\,\sigma_{2j+1}(n)
            - \frac{17}{150\,450\,048\,000}\,\tau(n),
\]
with explicit polynomials $P_j\in\Q[n]$:
\begin{align*}
P_0(n)&=\frac{550499}{4541644800}-\frac{153617}{371589120}\,n+\frac{2159}{6635520}\,n^2-\frac{67}{737280}\,n^3+\frac{11}{1105920}\,n^4-\frac{1}{2764800}\,n^5\\[4pt]
P_1(n)&=\frac{153617}{743178240}-\frac{2159}{8847360}\,n+\frac{67}{819200}\,n^2-\frac{11}{1105920}\,n^3+\frac{1}{2580480}\,n^4\\[4pt]
P_2(n)&=\frac{2159}{88473600}-\frac{67}{4915200}\,n+\frac{11}{5160960}\,n^2-\frac{1}{10321920}\,n^3\\[4pt]
P_3(n)&=\frac{67}{123863040}-\frac{11}{74317824}\,n+\frac{1}{111476736}\,n^2\\[4pt]
P_4(n)&=\frac{11}{3715891200}-\frac{1}{3096576000}\,n\\[4pt]
P_5(n)&=\frac{1}{268995133440}
\end{align*}
These 22 coefficients were verified against direct partition enumeration for
$n=1,\ldots,30$ with no mismatches (see repository script
\texttt{Paper/compute\_m6\_coefficients.jl}).
The key datum is the $\tau$-coefficient
\[
  c_\tau = -\frac{17}{150\,450\,048\,000},
  \qquad
  150\,450\,048\,000 = 2^{10}\times 3^5\times 5^3\times 7\times 691.
\]
The denominator prime $691 = \mathrm{num}(B_{12})$ is the Bernoulli prime
indexing the weight-12 Eisenstein series.

% ------------------------------------------------------------------
\subsection*{A.2\quad Cusp Structure of $M_7$ and Value of $\beta$}
\label{app:M7-formula}

Despite $\dim S_{14}(\SL_2(\Z))=0$, the closed form of $M_7(n)$ involves
$n\cdot\tau(n)$.
This follows from the quasimodular decomposition of $U_7(q)$: the depth-1 component
lies in $E_2\cdot M_{12}$, and the identity $D\Delta = E_2\Delta$ (proved from Ramanujan's differential equations
in the companion paper) shows that
$[q^n](E_2\Delta)=n\cdot\tau(n)$ exactly.
There is no independent $\tau(n)$ term: the $\tau(p)$ direction in $\mathcal{O}_d$
comes entirely from $M_6$ (via $c_\tau\ne 0$).

The exact closed form is
\[
  M_7(n) = \sum_{j=0}^{6} Q_j(n)\,\sigma_{2j+1}(n) + \beta\,n\cdot\tau(n),
\]
with $\beta\ne 0$.
The exact rational value of $\beta$ is the coefficient of $E_2\Delta$ in the
quasimodular decomposition of $U_7$, computed in exact arithmetic by the repository
script \texttt{Paper/extract\_e6\_e7.jl} (function \texttt{\_fit\_m7!}, column~30
of the augmented RREF system) and verified against direct partition enumeration for
$n=1,\dots,30$.

The structural significance is:
\begin{itemize}
  \item $\beta\ne 0$: $M_7$ contributes $p\cdot\tau(p)$ at primes, a direction absent
        from $M_6$ alone (which contributes only $\tau(p)$).
\end{itemize}
Together, the $\{M_6,M_7\}$ sub-system at primes carries both independent cusp
directions $\tau(p)$ and $p\cdot\tau(p)$, which is the source of the
cusp-cancellation mechanism underlying $E_6$.

% ------------------------------------------------------------------
\subsection*{A.3\quad Prime Evaluation Rank Data}
\label{app:rank-data}

\begin{table}[H]
\centering
\caption{Rank and nullity of prime evaluation matrices at $N=95$ primes,
         showing the pivot status of $M_6$ within the restricted basis.}
\label{tab:app-pivot}
\begin{tabular}{@{}lcccc@{}}
\toprule
Matrix & Rows & Cols & Rank & Nullity \\
\midrule
$a_{\max}=5,\;d=2$ & 95 & 15 & 11 & 4 \\
$a_{\max}=6,\;d=2$ & 95 & 18 & 14 & 4 \\
$a_{\max}=5,\;d=3$ & 95 & 20 & 12 & 8 \\
$a_{\max}=6,\;d=3$ & 95 & 24 & 16 & 8 \\
\bottomrule
\end{tabular}
\end{table}

In each case the nullity does not increase when $M_6$ is adjoined,
confirming that all $M_6$-columns are pivot columns.

% ------------------------------------------------------------------
\subsection*{A.4\quad Nullity Gains When $M_7$ Is Adjoined}
\label{app:nullity-m7}

\begin{table}[H]
\centering
\caption{Nullity comparison at $N=95$ primes.
         The gain first appears at $d=4$.}
\label{tab:app-nullity-m7}
\begin{tabular}{@{}ccccc@{}}
\toprule
$d$ & Nullity $\{M_1\text{--}M_5\}$ &
      Nullity $\{M_1\text{--}M_6\}$ &
      Nullity $\{M_1\text{--}M_7\}$ & Gain \\
\midrule
2 & 4  & 4  & 4  & 0 \\
3 & 8  & 8  & 8  & 0 \\
4 & 12 & 12 & 13 & $+1$ \\
5 & 16 & 16 & 18 & $+2$ \\
6 & 20 & 20 & 23 & $+3$ \\
\bottomrule
\end{tabular}
\end{table}

% ------------------------------------------------------------------
\subsection*{A.5\quad Explicit Formula for $E_5$}
\label{app:E5-formula}

After normalization to primitive integer coefficients,
the unique new prime-vanishing direction at degree~2 in the basis
$\{M_1,\dots,M_5\}$ is
\begin{align*}
E_5(n) ={}
  & (-450450+675675n-225225n^2)\,M_1(n) \\
  &+(960960n-120120n^2)\,M_2(n) \\
  &+(2534912n-166016n^2)\,M_3(n) \\
  &+(7999488n-322560n^2)\,M_4(n) \\
  &+258048000\,M_5(n).
\end{align*}
The constant $M_5$-coefficient is
$258048000 = 2^{11}\times 3^2\times 5^3\times 7\times 127$.

% ------------------------------------------------------------------
\subsection*{A.6\quad Explicit Formula for $E_6$}
\label{app:E6-formula}

The unique new prime-vanishing direction at degree~4 in the basis
$\{M_1,\dots,M_7\}$, normalized analogously, is
\begin{align*}
E_6(n) ={}
  &(105367732470 - 277943208789n + 289613157079n^2 \\
  &\quad{}-145583618619n^3+28545937859n^4)\,M_1(n)\\
  &+(-51324522800n^3+4292382480n^4)\,M_2(n)\\
  &+(-40313554176n^3+1776519808n^4)\,M_3(n)\\
  &+(-32888346624n^3+770273280n^4)\,M_4(n)\\
  &+(-7741440000n^3-154828800n^4)\,M_5(n)\\
  &+(-483548921856000+37196070912000n)\,M_6(n)\\
  &+892705701888000\,M_7(n).
\end{align*}
The leading $M_7$-coefficient is
$892705701888000 = 2^{10}\times 3^5\times 5^3\times 7\times 691\times 1303$.
The factor $691$ connects $E_6$ arithmetically to the weight-12 modular structure of
$M_6$.

% ------------------------------------------------------------------
\subsection*{A.7\quad Extended Search to $a_{\max}=12$}
\label{app:extended-search}

\begin{table}[H]
\centering
\caption{Extended search to $a_{\max}=12$.
         No independent $E_i$ was found beyond $E_6$.}
\label{tab:app-extended}
\begin{tabular}{@{}ccccc@{}}
\toprule
$a$ & Weight $2a$ & $\dim S_{2a}$ & New cusp type & New $E_i$ found? \\
\midrule
8  & 16 & 1 & $\Delta\cdot E_4$              & No \\
9  & 18 & 1 & $\Delta\cdot E_6$              & No \\
10 & 20 & 1 & $\Delta\cdot E_8$              & No \\
11 & 22 & 1 & $\Delta\cdot E_{10}$           & No \\
12 & 24 & 2 & $\Delta\cdot E_{12},\;\Delta^2$ & No \\
\bottomrule
\end{tabular}
\end{table}


\section*{Appendix: Computational Data}
\addcontentsline{toc}{section}{Appendix: Computational Data}
\label{app:computations}

All computations were performed in exact rational arithmetic.
Prime evaluation matrices were built from the first $N$ primes and reduced via
rational row-reduction (RREF over $\Q$).
The source code is available at the project repository.

% ------------------------------------------------------------------
\subsection*{A.1\quad Closed Form of $M_6$}
\label{app:M6-formula}

Fitting the $55\times 22$ rational linear system
(21 polynomial-weighted divisor-sum columns plus one $\tau$-column)
yields the exact formula
\[
  M_6(n) = \sum_{j=0}^{5} P_j(n)\,\sigma_{2j+1}(n)
            - \frac{17}{150\,450\,048\,000}\,\tau(n),
\]
with explicit polynomials $P_j\in\Q[n]$:
\begin{align*}
P_0(n)&=\frac{550499}{4541644800}-\frac{153617}{371589120}\,n+\frac{2159}{6635520}\,n^2-\frac{67}{737280}\,n^3+\frac{11}{1105920}\,n^4-\frac{1}{2764800}\,n^5\\[4pt]
P_1(n)&=\frac{153617}{743178240}-\frac{2159}{8847360}\,n+\frac{67}{819200}\,n^2-\frac{11}{1105920}\,n^3+\frac{1}{2580480}\,n^4\\[4pt]
P_2(n)&=\frac{2159}{88473600}-\frac{67}{4915200}\,n+\frac{11}{5160960}\,n^2-\frac{1}{10321920}\,n^3\\[4pt]
P_3(n)&=\frac{67}{123863040}-\frac{11}{74317824}\,n+\frac{1}{111476736}\,n^2\\[4pt]
P_4(n)&=\frac{11}{3715891200}-\frac{1}{3096576000}\,n\\[4pt]
P_5(n)&=\frac{1}{268995133440}
\end{align*}
These 22 coefficients were verified against direct partition enumeration for
$n=1,\ldots,30$ with no mismatches (see repository script
\texttt{Paper/compute\_m6\_coefficients.jl}).
The key datum is the $\tau$-coefficient
\[
  c_\tau = -\frac{17}{150\,450\,048\,000},
  \qquad
  150\,450\,048\,000 = 2^{10}\times 3^5\times 5^3\times 7\times 691.
\]
The denominator prime $691 = \mathrm{num}(B_{12})$ is the Bernoulli prime
indexing the weight-12 Eisenstein series.

% ------------------------------------------------------------------
\subsection*{A.2\quad Cusp Structure of $M_7$ and Value of $\beta$}
\label{app:M7-formula}

Despite $\dim S_{14}(\SL_2(\Z))=0$, the closed form of $M_7(n)$ involves
$n\cdot\tau(n)$.
This follows from the quasimodular decomposition of $U_7(q)$: the depth-1 component
lies in $E_2\cdot M_{12}$, and the identity $D\Delta = E_2\Delta$ (proved from Ramanujan's differential equations
in the companion paper) shows that
$[q^n](E_2\Delta)=n\cdot\tau(n)$ exactly.
There is no independent $\tau(n)$ term: the $\tau(p)$ direction in $\mathcal{O}_d$
comes entirely from $M_6$ (via $c_\tau\ne 0$).

The exact closed form is
\[
  M_7(n) = \sum_{j=0}^{6} Q_j(n)\,\sigma_{2j+1}(n) + \beta\,n\cdot\tau(n),
\]
with $\beta\ne 0$.
The exact rational value of $\beta$ is the coefficient of $E_2\Delta$ in the
quasimodular decomposition of $U_7$, computed in exact arithmetic by the repository
script \texttt{Paper/extract\_e6\_e7.jl} (function \texttt{\_fit\_m7!}, column~30
of the augmented RREF system) and verified against direct partition enumeration for
$n=1,\dots,30$.

The structural significance is:
\begin{itemize}
  \item $\beta\ne 0$: $M_7$ contributes $p\cdot\tau(p)$ at primes, a direction absent
        from $M_6$ alone (which contributes only $\tau(p)$).
\end{itemize}
Together, the $\{M_6,M_7\}$ sub-system at primes carries both independent cusp
directions $\tau(p)$ and $p\cdot\tau(p)$, which is the source of the
cusp-cancellation mechanism underlying $E_6$.

% ------------------------------------------------------------------
\subsection*{A.3\quad Prime Evaluation Rank Data}
\label{app:rank-data}

\begin{table}[H]
\centering
\caption{Rank and nullity of prime evaluation matrices at $N=95$ primes,
         showing the pivot status of $M_6$ within the restricted basis.}
\label{tab:app-pivot}
\begin{tabular}{@{}lcccc@{}}
\toprule
Matrix & Rows & Cols & Rank & Nullity \\
\midrule
$a_{\max}=5,\;d=2$ & 95 & 15 & 11 & 4 \\
$a_{\max}=6,\;d=2$ & 95 & 18 & 14 & 4 \\
$a_{\max}=5,\;d=3$ & 95 & 20 & 12 & 8 \\
$a_{\max}=6,\;d=3$ & 95 & 24 & 16 & 8 \\
\bottomrule
\end{tabular}
\end{table}

In each case the nullity does not increase when $M_6$ is adjoined,
confirming that all $M_6$-columns are pivot columns.

% ------------------------------------------------------------------
\subsection*{A.4\quad Nullity Gains When $M_7$ Is Adjoined}
\label{app:nullity-m7}

\begin{table}[H]
\centering
\caption{Nullity comparison at $N=95$ primes.
         The gain first appears at $d=4$.}
\label{tab:app-nullity-m7}
\begin{tabular}{@{}ccccc@{}}
\toprule
$d$ & Nullity $\{M_1\text{--}M_5\}$ &
      Nullity $\{M_1\text{--}M_6\}$ &
      Nullity $\{M_1\text{--}M_7\}$ & Gain \\
\midrule
2 & 4  & 4  & 4  & 0 \\
3 & 8  & 8  & 8  & 0 \\
4 & 12 & 12 & 13 & $+1$ \\
5 & 16 & 16 & 18 & $+2$ \\
6 & 20 & 20 & 23 & $+3$ \\
\bottomrule
\end{tabular}
\end{table}

% ------------------------------------------------------------------
\subsection*{A.5\quad Explicit Formula for $E_5$}
\label{app:E5-formula}

After normalization to primitive integer coefficients,
the unique new prime-vanishing direction at degree~2 in the basis
$\{M_1,\dots,M_5\}$ is
\begin{align*}
E_5(n) ={}
  & (-450450+675675n-225225n^2)\,M_1(n) \\
  &+(960960n-120120n^2)\,M_2(n) \\
  &+(2534912n-166016n^2)\,M_3(n) \\
  &+(7999488n-322560n^2)\,M_4(n) \\
  &+258048000\,M_5(n).
\end{align*}
The constant $M_5$-coefficient is
$258048000 = 2^{11}\times 3^2\times 5^3\times 7\times 127$.

% ------------------------------------------------------------------
\subsection*{A.6\quad Explicit Formula for $E_6$}
\label{app:E6-formula}

The unique new prime-vanishing direction at degree~4 in the basis
$\{M_1,\dots,M_7\}$, normalized analogously, is
\begin{align*}
E_6(n) ={}
  &(105367732470 - 277943208789n + 289613157079n^2 \\
  &\quad{}-145583618619n^3+28545937859n^4)\,M_1(n)\\
  &+(-51324522800n^3+4292382480n^4)\,M_2(n)\\
  &+(-40313554176n^3+1776519808n^4)\,M_3(n)\\
  &+(-32888346624n^3+770273280n^4)\,M_4(n)\\
  &+(-7741440000n^3-154828800n^4)\,M_5(n)\\
  &+(-483548921856000+37196070912000n)\,M_6(n)\\
  &+892705701888000\,M_7(n).
\end{align*}
The leading $M_7$-coefficient is
$892705701888000 = 2^{10}\times 3^5\times 5^3\times 7\times 691\times 1303$.
The factor $691$ connects $E_6$ arithmetically to the weight-12 modular structure of
$M_6$.

% ------------------------------------------------------------------
\subsection*{A.7\quad Extended Search to $a_{\max}=12$}
\label{app:extended-search}

\begin{table}[H]
\centering
\caption{Extended search to $a_{\max}=12$.
         No independent $E_i$ was found beyond $E_6$.}
\label{tab:app-extended}
\begin{tabular}{@{}ccccc@{}}
\toprule
$a$ & Weight $2a$ & $\dim S_{2a}$ & New cusp type & New $E_i$ found? \\
\midrule
8  & 16 & 1 & $\Delta\cdot E_4$              & No \\
9  & 18 & 1 & $\Delta\cdot E_6$              & No \\
10 & 20 & 1 & $\Delta\cdot E_8$              & No \\
11 & 22 & 1 & $\Delta\cdot E_{10}$           & No \\
12 & 24 & 2 & $\Delta\cdot E_{12},\;\Delta^2$ & No \\
\bottomrule
\end{tabular}
\end{table}


\section*{Appendix: Computational Data}
\addcontentsline{toc}{section}{Appendix: Computational Data}
\label{app:computations}

All computations were performed in exact rational arithmetic.
Prime evaluation matrices were built from the first $N$ primes and reduced via
rational row-reduction (RREF over $\Q$).
The source code is available at the project repository.

% ------------------------------------------------------------------
\subsection*{A.1\quad Closed Form of $M_6$}
\label{app:M6-formula}

Fitting the $55\times 22$ rational linear system
(21 polynomial-weighted divisor-sum columns plus one $\tau$-column)
yields the exact formula
\[
  M_6(n) = \sum_{j=0}^{5} P_j(n)\,\sigma_{2j+1}(n)
            - \frac{17}{150\,450\,048\,000}\,\tau(n),
\]
with explicit polynomials $P_j\in\Q[n]$:
\begin{align*}
P_0(n)&=\frac{550499}{4541644800}-\frac{153617}{371589120}\,n+\frac{2159}{6635520}\,n^2-\frac{67}{737280}\,n^3+\frac{11}{1105920}\,n^4-\frac{1}{2764800}\,n^5\\[4pt]
P_1(n)&=\frac{153617}{743178240}-\frac{2159}{8847360}\,n+\frac{67}{819200}\,n^2-\frac{11}{1105920}\,n^3+\frac{1}{2580480}\,n^4\\[4pt]
P_2(n)&=\frac{2159}{88473600}-\frac{67}{4915200}\,n+\frac{11}{5160960}\,n^2-\frac{1}{10321920}\,n^3\\[4pt]
P_3(n)&=\frac{67}{123863040}-\frac{11}{74317824}\,n+\frac{1}{111476736}\,n^2\\[4pt]
P_4(n)&=\frac{11}{3715891200}-\frac{1}{3096576000}\,n\\[4pt]
P_5(n)&=\frac{1}{268995133440}
\end{align*}
These 22 coefficients were verified against direct partition enumeration for
$n=1,\ldots,30$ with no mismatches (see repository script
\texttt{Paper/compute\_m6\_coefficients.jl}).
The key datum is the $\tau$-coefficient
\[
  c_\tau = -\frac{17}{150\,450\,048\,000},
  \qquad
  150\,450\,048\,000 = 2^{10}\times 3^5\times 5^3\times 7\times 691.
\]
The denominator prime $691 = \mathrm{num}(B_{12})$ is the Bernoulli prime
indexing the weight-12 Eisenstein series.

% ------------------------------------------------------------------
\subsection*{A.2\quad Cusp Structure of $M_7$ and Value of $\beta$}
\label{app:M7-formula}

Despite $\dim S_{14}(\SL_2(\Z))=0$, the closed form of $M_7(n)$ involves
$n\cdot\tau(n)$.
This follows from the quasimodular decomposition of $U_7(q)$: the depth-1 component
lies in $E_2\cdot M_{12}$, and the identity $D\Delta = E_2\Delta$ (proved from Ramanujan's differential equations
in the companion paper) shows that
$[q^n](E_2\Delta)=n\cdot\tau(n)$ exactly.
There is no independent $\tau(n)$ term: the $\tau(p)$ direction in $\mathcal{O}_d$
comes entirely from $M_6$ (via $c_\tau\ne 0$).

The exact closed form is
\[
  M_7(n) = \sum_{j=0}^{6} Q_j(n)\,\sigma_{2j+1}(n) + \beta\,n\cdot\tau(n),
\]
with $\beta\ne 0$.
The exact rational value of $\beta$ is the coefficient of $E_2\Delta$ in the
quasimodular decomposition of $U_7$, computed in exact arithmetic by the repository
script \texttt{Paper/extract\_e6\_e7.jl} (function \texttt{\_fit\_m7!}, column~30
of the augmented RREF system) and verified against direct partition enumeration for
$n=1,\dots,30$.

The structural significance is:
\begin{itemize}
  \item $\beta\ne 0$: $M_7$ contributes $p\cdot\tau(p)$ at primes, a direction absent
        from $M_6$ alone (which contributes only $\tau(p)$).
\end{itemize}
Together, the $\{M_6,M_7\}$ sub-system at primes carries both independent cusp
directions $\tau(p)$ and $p\cdot\tau(p)$, which is the source of the
cusp-cancellation mechanism underlying $E_6$.

% ------------------------------------------------------------------
\subsection*{A.3\quad Prime Evaluation Rank Data}
\label{app:rank-data}

\begin{table}[H]
\centering
\caption{Rank and nullity of prime evaluation matrices at $N=95$ primes,
         showing the pivot status of $M_6$ within the restricted basis.}
\label{tab:app-pivot}
\begin{tabular}{@{}lcccc@{}}
\toprule
Matrix & Rows & Cols & Rank & Nullity \\
\midrule
$a_{\max}=5,\;d=2$ & 95 & 15 & 11 & 4 \\
$a_{\max}=6,\;d=2$ & 95 & 18 & 14 & 4 \\
$a_{\max}=5,\;d=3$ & 95 & 20 & 12 & 8 \\
$a_{\max}=6,\;d=3$ & 95 & 24 & 16 & 8 \\
\bottomrule
\end{tabular}
\end{table}

In each case the nullity does not increase when $M_6$ is adjoined,
confirming that all $M_6$-columns are pivot columns.

% ------------------------------------------------------------------
\subsection*{A.4\quad Nullity Gains When $M_7$ Is Adjoined}
\label{app:nullity-m7}

\begin{table}[H]
\centering
\caption{Nullity comparison at $N=95$ primes.
         The gain first appears at $d=4$.}
\label{tab:app-nullity-m7}
\begin{tabular}{@{}ccccc@{}}
\toprule
$d$ & Nullity $\{M_1\text{--}M_5\}$ &
      Nullity $\{M_1\text{--}M_6\}$ &
      Nullity $\{M_1\text{--}M_7\}$ & Gain \\
\midrule
2 & 4  & 4  & 4  & 0 \\
3 & 8  & 8  & 8  & 0 \\
4 & 12 & 12 & 13 & $+1$ \\
5 & 16 & 16 & 18 & $+2$ \\
6 & 20 & 20 & 23 & $+3$ \\
\bottomrule
\end{tabular}
\end{table}

% ------------------------------------------------------------------
\subsection*{A.5\quad Explicit Formula for $E_5$}
\label{app:E5-formula}

After normalization to primitive integer coefficients,
the unique new prime-vanishing direction at degree~2 in the basis
$\{M_1,\dots,M_5\}$ is
\begin{align*}
E_5(n) ={}
  & (-450450+675675n-225225n^2)\,M_1(n) \\
  &+(960960n-120120n^2)\,M_2(n) \\
  &+(2534912n-166016n^2)\,M_3(n) \\
  &+(7999488n-322560n^2)\,M_4(n) \\
  &+258048000\,M_5(n).
\end{align*}
The constant $M_5$-coefficient is
$258048000 = 2^{11}\times 3^2\times 5^3\times 7\times 127$.

% ------------------------------------------------------------------
\subsection*{A.6\quad Explicit Formula for $E_6$}
\label{app:E6-formula}

The unique new prime-vanishing direction at degree~4 in the basis
$\{M_1,\dots,M_7\}$, normalized analogously, is
\begin{align*}
E_6(n) ={}
  &(105367732470 - 277943208789n + 289613157079n^2 \\
  &\quad{}-145583618619n^3+28545937859n^4)\,M_1(n)\\
  &+(-51324522800n^3+4292382480n^4)\,M_2(n)\\
  &+(-40313554176n^3+1776519808n^4)\,M_3(n)\\
  &+(-32888346624n^3+770273280n^4)\,M_4(n)\\
  &+(-7741440000n^3-154828800n^4)\,M_5(n)\\
  &+(-483548921856000+37196070912000n)\,M_6(n)\\
  &+892705701888000\,M_7(n).
\end{align*}
The leading $M_7$-coefficient is
$892705701888000 = 2^{10}\times 3^5\times 5^3\times 7\times 691\times 1303$.
The factor $691$ connects $E_6$ arithmetically to the weight-12 modular structure of
$M_6$.

% ------------------------------------------------------------------
\subsection*{A.7\quad Extended Search to $a_{\max}=12$}
\label{app:extended-search}

\begin{table}[H]
\centering
\caption{Extended search to $a_{\max}=12$.
         No independent $E_i$ was found beyond $E_6$.}
\label{tab:app-extended}
\begin{tabular}{@{}ccccc@{}}
\toprule
$a$ & Weight $2a$ & $\dim S_{2a}$ & New cusp type & New $E_i$ found? \\
\midrule
8  & 16 & 1 & $\Delta\cdot E_4$              & No \\
9  & 18 & 1 & $\Delta\cdot E_6$              & No \\
10 & 20 & 1 & $\Delta\cdot E_8$              & No \\
11 & 22 & 1 & $\Delta\cdot E_{10}$           & No \\
12 & 24 & 2 & $\Delta\cdot E_{12},\;\Delta^2$ & No \\
\bottomrule
\end{tabular}
\end{table}


\end{document}
